%% ----------------------------------------------------------------------------
%%  §6.9  Esercizi di riepilogo
%% ----------------------------------------------------------------------------
\section{Esercizi di riepilogo}
\label{sec:esercizi-cap6}

Come nei capitoli precedenti, ogni soluzione è accompagnata da una nota sulla
\emph{scelta logica} che la guida.

\needspace{8\baselineskip}
\subsection*{I teoremi di Rolle, di Lagrange e di Darboux}

\begin{esercizio}\label{es:cubica-unica}
Dimostrare che l'equazione $x^{3}+x-1=0$ \eqref{eq:cubica} ha una e una sola soluzione
reale, e che essa appartiene all'intervallo $(0,1)$.
\begin{soluzione}
Sia $f(x)=x^{3}+x-1$. \emph{Esistenza:} $f$ è continua, $f(0)=-1<0$ e $f(1)=1>0$; per il
teorema degli zeri (\cref{teo:zeri}) esiste una soluzione in $(0,1)$. \emph{Unicità:} se
ci fossero due soluzioni $x_{1}<x_{2}$, sarebbe $f(x_{1})=f(x_{2})=0$ e, per il teorema di
Rolle, $f'(\xi)=3\xi^{2}+1=0$ per qualche $\xi$, il che è impossibile perché
$3\xi^{2}+1\ge1$. In alternativa: $f'>0$ su $\R$, quindi $f$ è strettamente crescente
(\cref{cor:derivata-positiva}) e iniettiva. \emph{Scelta logica:} esistenza e unicità
richiedono strumenti diversi; il teorema degli zeri non dice nulla sul numero delle
soluzioni, e il teorema di Rolle (o la monotonia) lo controlla: tra due zeri di $f$ c'è
uno zero di $f'$.
\end{soluzione}
\end{esercizio}

\begin{esercizio}\label{es:log-lagrange}
Dimostrare con il teorema di Lagrange che, per ogni $x>0$,
$\dfrac{x}{1+x}<\log(1+x)<x$.
\begin{soluzione}
Applichiamo il teorema di Lagrange a $\log t$ nell'intervallo $[1,1+x]$: esiste
$\xi\in(1,1+x)$ tale che $\log(1+x)-\log1=\frac1\xi\,x$. Poiché $1<\xi<1+x$, risulta
$\frac1{1+x}<\frac1\xi<1$, e moltiplicando per $x>0$ si ottiene la tesi. La
disuguaglianza di destra era già nell'\cref{es:disuguaglianze-exp-log}.
\emph{Scelta logica:} il teorema di Lagrange trasforma un'uguaglianza esatta ma con un
punto sconosciuto in una doppia disuguaglianza, sostituendo al punto sconosciuto gli
estremi dell'intervallo in cui si trova; è il meccanismo dell'\cref{es:stima-radice}.
\end{soluzione}
\end{esercizio}

\begin{esercizio}\label{es:darboux}
(\emph{Teorema di Darboux}.) Sia $g$ derivabile in $[a,b]$ e sia $\lambda$ un numero
strettamente compreso tra $g'(a)$ e $g'(b)$. Dimostrare che esiste $\xi\in(a,b)$ con
$g'(\xi)=\lambda$. Dedurne che la funzione che vale $0$ per $x<0$ e $1$ per $x\ge0$ non è
la derivata di alcuna funzione in $[-1,1]$.
\begin{soluzione}
La funzione $\varphi(x)=g(x)-\lambda x$ è derivabile in $[a,b]$, con
$\varphi'(x)=g'(x)-\lambda$; per la scelta di $\lambda$, $\varphi'(a)$ e $\varphi'(b)$ hanno
segni opposti. Se $\varphi'$ non si annullasse in alcun punto di $[a,b]$, per il
\cref{lem:darboux} avrebbe segno costante: contraddizione. Dunque $\varphi'(\xi)=0$ per
qualche $\xi$, che non può essere né $a$ né $b$ (dove $\varphi'\neq0$); cioè
$g'(\xi)=\lambda$ con $\xi\in(a,b)$. Una funzione con un salto non assume i valori
intermedi, e quindi non può essere una derivata. \emph{Scelta logica:} si sposta il
problema «$g'$ assume il valore $\lambda$» nel problema «$\varphi'$ si annulla»,
sottraendo a $g$ la retta di pendenza $\lambda$; è lo stesso trucco della dimostrazione
del teorema di Lagrange.
\end{soluzione}
\end{esercizio}

\begin{esercizio}\label{es:corde-convesse}
Sia $f$ convessa in $[a,b]$ nel senso della \cref{def:convessa}. Dimostrare che ogni
corda sta sopra il grafico, cioè che
$f\bigl((1-t)x+ty\bigr)\le(1-t)f(x)+tf(y)$ per ogni $x,y\in[a,b]$ e ogni $t\in[0,1]$.
\begin{soluzione}
Sia $z=(1-t)x+ty\in[a,b]$. Scriviamo la \eqref{eq:convessa} con punto di tangenza $z$,
prima nel punto $x$ e poi nel punto $y$:
$f(x)\ge f(z)+f'(z)(x-z)$, $f(y)\ge f(z)+f'(z)(y-z)$. Moltiplichiamo la prima per
$1-t\ge0$, la seconda per $t\ge0$ e sommiamo:
\[
  (1-t)f(x)+tf(y)\ge f(z)+f'(z)\bigl[(1-t)x+ty-z\bigr]=f(z) .
\]
\emph{Scelta logica:} si sceglie come punto di tangenza proprio il punto $z$ in cui si
vuole stimare $f$; la combinazione con pesi $1-t$ e $t$ è quella che fa sparire il termine
con la derivata, perché $(1-t)x+ty-z=0$.
\end{soluzione}
\end{esercizio}

\needspace{8\baselineskip}
\subsection*{Disuguaglianze e problemi di ottimo}

\begin{esercizio}\label{es:bernoulli-convessita}
Dimostrare che $(1+x)^{b}\ge1+bx$ per ogni $x>-1$ e ogni $b\ge1$, e che la disuguaglianza
si inverte se $0\le b\le1$. Confrontare con la disuguaglianza di Bernoulli
(\cref{teo:bernoulli}).
\begin{soluzione}
La funzione $f(x)=(1+x)^{b}$ è derivabile due volte in $(-1,+\infty)$, con
$f''(x)=b(b-1)(1+x)^{b-2}$, che è $\ge0$ se $b\ge1$ e $\le0$ se $0\le b\le1$. Nel primo
caso $f$ è convessa, e la \eqref{eq:convessa} con $x_{0}=0$ dà
$(1+x)^{b}\ge f(0)+f'(0)x=1+bx$; nel secondo caso $f$ è concava e la disuguaglianza si
inverte. La disuguaglianza di Bernoulli del \cref{cap:numeri}, dimostrata per induzione
per gli esponenti naturali, si estende così a ogni esponente reale $b\ge1$.
\emph{Scelta logica:} $y=1+bx$ è la retta tangente al grafico di $(1+x)^{b}$ nel punto
$x=0$, e la disuguaglianza dice da che parte della tangente sta il grafico.
\end{soluzione}
\end{esercizio}

\begin{esercizio}[La trave di massima resistenza]\label{es:trave}
Da un tronco cilindrico di diametro $d$ si vuole ricavare una trave a sezione
rettangolare, di base $b$ e altezza $h$, inflessa nel piano dell'altezza. La resistenza a
flessione è proporzionale al \emph{modulo di resistenza} $W_{f}=\frac{bh^{2}}6$.
Determinare le dimensioni che rendono massimo $W_{f}$.
\begin{soluzione}
Il rettangolo è inscritto nella circonferenza, quindi $b^{2}+h^{2}=d^{2}$ e
\[
  W_{f}(b)=\frac{b(d^{2}-b^{2})}{6}, \qquad b\in[0,d] .
\]
La funzione è continua nell'intervallo chiuso, nulla negli estremi e positiva
all'interno; per il metodo del \cref{sec:fermat} il massimo assoluto è in un punto
stazionario interno. Da $W_{f}'(b)=\frac{d^{2}-3b^{2}}6=0$ si ricava $b^{*}=\frac d{\sqrt3}$,
e quindi $h^{*}=\sqrt{d^{2}-b^{*2}}=d\sqrt{\frac23}$: il rapporto ottimo è
$\frac{h^{*}}{b^{*}}=\sqrt2$. Inoltre $W_{f}''(b)=-b<0$ conferma il massimo. Se invece
si vuole la trave più \emph{rigida}, si massimizza il momento d'inerzia del
second'ordine $\frac{bh^{3}}{12}$, e si trova $\frac hb=\sqrt3$. Nelle costruzioni
aeronautiche lo stesso tipo di calcolo guida il dimensionamento delle solette dei
longheroni. \emph{Scelta logica:} il vincolo geometrico permette di eliminare una
variabile e di ridurre il problema a una funzione di una sola variabile su un intervallo
chiuso; il teorema di Weierstrass garantisce l'esistenza del massimo, e il teorema di
Fermat lo localizza.
\end{soluzione}
\end{esercizio}

\begin{python}[La trave di massima resistenza con sympy]
import sympy as sp

b, d = sp.symbols("b d", positive=True)
h2 = d**2 - b**2                  # Pitagora: b^2 + h^2 = d^2
Wf = b * h2 / 6                   # modulo di resistenza b h^2 / 6
bs = sp.solve(sp.diff(Wf, b), b)[0]
print("b* =", bs, "  h* =", sp.sqrt(h2.subs(b, bs)),
      "  h*/b* =", sp.simplify(sp.sqrt(h2.subs(b, bs)) / bs))
print("W'' (b*) =", sp.simplify(sp.diff(Wf, b, 2).subs(b, bs)))
print("W_max =", sp.simplify(Wf.subs(b, bs)))
\end{python}
\begin{risultato}
b* = sqrt(3)*d/3   h* = sqrt(6)*d/3   h*/b* = sqrt(2)
W'' (b*) = -sqrt(3)*d/3
W_max = sqrt(3)*d**3/27
\end{risultato}

\begin{esercizio}[Il principio di Fermat e la legge di Snell]\label{es:snell}
Un raggio di luce va dal punto $A=(0,h_{1})$, nel mezzo $y>0$ dove la velocità della
luce è $v_{1}$, al punto $B=(L,-h_{2})$, nel mezzo $y<0$ dove la velocità è $v_{2}$,
attraversando l'asse $x$ nel punto $(x,0)$. Dimostrare che il tempo di percorrenza è
minimo quando
\begin{equation}\label{eq:snell}
  \frac{\sen\theta_{1}}{v_{1}}=\frac{\sen\theta_{2}}{v_{2}},
\end{equation}
dove $\theta_{1}$, $\theta_{2}$ sono gli angoli che i due tratti del percorso formano con
la normale alla superficie di separazione.
\begin{soluzione}
Il tempo di percorrenza è
\begin{align*}
  T(x)&=\frac{\sqrt{h_{1}^{2}+x^{2}}}{v_{1}}+\frac{\sqrt{h_{2}^{2}+(L-x)^{2}}}{v_{2}},\\
  T'(x)&=\frac{x}{v_{1}\sqrt{h_{1}^{2}+x^{2}}}-\frac{L-x}{v_{2}\sqrt{h_{2}^{2}+(L-x)^{2}}} .
\end{align*}
Le frazioni $\frac{x}{\sqrt{h_{1}^{2}+x^{2}}}$ e $\frac{L-x}{\sqrt{h_{2}^{2}+(L-x)^{2}}}$ sono
proprio $\sen\theta_{1}$ e $\sen\theta_{2}$; quindi $T'(x)=0$ equivale alla
\eqref{eq:snell}. Resta da provare che il punto stazionario esiste, è unico ed è il
minimo. Si ha $T'(0)<0<T'(L)$, e un calcolo diretto dà
\[
  T''(x)=\frac{h_{1}^{2}}{v_{1}(h_{1}^{2}+x^{2})^{3/2}}+\frac{h_{2}^{2}}{v_{2}(h_{2}^{2}+(L-x)^{2})^{3/2}}>0 ,
\]
quindi $T'$ è strettamente crescente: per il teorema degli zeri si annulla in un unico
punto $x^{*}\in(0,L)$, nel quale $T'$ passa da negativa a positiva. Per il criterio di
monotonia, $T$ è strettamente decrescente in $[0,x^{*}]$ e strettamente crescente in
$[x^{*},L]$, e quindi $x^{*}$ è il punto di minimo assoluto di $T$ in $[0,L]$ (e anche su $\R$,
perché $T$ è strettamente convessa). Il programma seguente risolve numericamente
l'equazione $T'(x)=0$ per $h_{1}=h_{2}=1$, $L=2$, $v_{1}=1$, $v_{2}=0{,}75$ e verifica la
\eqref{eq:snell}. \emph{Scelta logica:} la convessità risolve in un colpo solo
esistenza, unicità e natura del punto stazionario.
\end{soluzione}
\end{esercizio}

\begin{python}[Il percorso di tempo minimo]
import math
from scipy.optimize import brentq

# raggio da A=(0, h1) a B=(L, -h2); velocità v1 sopra, v2 sotto
h1, h2, L, v1, v2 = 1.0, 1.0, 2.0, 1.0, 0.75
T = lambda x: math.hypot(h1, x) / v1 + math.hypot(h2, L - x) / v2
dT = lambda x: (x / math.hypot(h1, x) / v1
                - (L - x) / math.hypot(h2, L - x) / v2)
xs = brentq(dT, 0, L)                    # T'(x) = 0, unico zero
s1, s2 = xs / math.hypot(h1, xs), (L - xs) / math.hypot(h2, L - xs)
print(f"x* = {xs:.6f}   T(x*) = {T(xs):.6f}")
print(f"sen t1 / v1 = {s1 / v1:.6f}   sen t2 / v2 = {s2 / v2:.6f}")
print(f"T(0) = {T(0):.6f}   T(L) = {T(L):.6f}")
\end{python}
\begin{risultato}
x* = 1.270505   T(x*) = 3.267251
sen t1 / v1 = 0.785794   sen t2 / v2 = 0.785794
T(0) = 3.981424   T(L) = 3.569401
\end{risultato}

\needspace{8\baselineskip}
\subsection*{Limiti con il teorema di L'Hôpital e con la formula di Taylor}

\begin{esercizio}
Calcolare $\displaystyle\lim_{x\to0}\frac{e^{x}-e^{-x}-2x}{x-\sen x}$ in due modi: con il
teorema di L'Hôpital e con la formula di Taylor.
\begin{soluzione}
È una forma $\frac00$. Derivando numeratore e denominatore si ottengono
successivamente i rapporti
\[
  \frac{e^{x}+e^{-x}-2}{1-\cos x}, \qquad \frac{e^{x}-e^{-x}}{\sen x}, \qquad
  \frac{e^{x}+e^{-x}}{\cos x};
\]
i primi due sono ancora forme $\frac00$, l'ultimo tende a $2$. Applicando tre volte il
teorema di L'Hôpital, e risalendo la catena, il limite cercato vale $2$. Con la formula
di Taylor \eqref{eq:taylor-exp} di ordine $3$, applicata a $e^{x}$ e a $e^{-x}$,
\[
\begin{aligned}
  e^{x}-e^{-x}-2x&=\Bigl(1+x+\frac{x^{2}}2+\frac{x^{3}}6\Bigr)
  -\Bigl(1-x+\frac{x^{2}}2-\frac{x^{3}}6\Bigr)-2x+o(x^{3})\\
  &=\frac{x^{3}}3+o(x^{3})
\end{aligned}
\]
(i termini di grado pari si cancellano), e $x-\sen x=\frac{x^{3}}6+o(x^{4})$ per la
\eqref{eq:taylor-sen}; quindi il rapporto tende a $\frac{1/3}{1/6}=2$.
\emph{Scelta logica:} la formula di Taylor rivela subito che numeratore e denominatore
sono infinitesimi del terzo ordine, il che spiega perché il teorema di L'Hôpital
richiede tre applicazioni.
\end{soluzione}
\end{esercizio}

\needspace{8\baselineskip}
\begin{esercizio}
Calcolare $\displaystyle\lim_{x\to+\infty}x\Bigl(\frac\pi2-\operatorname{arctg}x\Bigr)$.
\begin{soluzione}
È una forma $\infty\cdot0$. Scrivendola come quoziente,
\[
  \lim_{x\to+\infty}\frac{\frac\pi2-\operatorname{arctg}x}{1/x}
  =\lim_{x\to+\infty}\frac{-\frac1{1+x^{2}}}{-\frac1{x^{2}}}
  =\lim_{x\to+\infty}\frac{x^{2}}{1+x^{2}}=1 .
\]
Il risultato dice che $\operatorname{arctg}x=\frac\pi2-\frac1x+o\bigl(\frac1x\bigr)$ per
$x\to+\infty$: il grafico si avvicina all'asintoto $y=\frac\pi2$ come l'iperbole
$y=\frac\pi2-\frac1x$. \emph{Scelta logica:} conviene mettere a denominatore il fattore
$\frac1x$, la cui derivata è semplice, e lasciare a numeratore la funzione trascendente,
la cui derivata è razionale.
\end{soluzione}
\end{esercizio}

\begin{esercizio}
Stabilire la natura del punto stazionario $x=0$ della funzione
$f(x)=e^{x}+e^{-x}-x^{2}-2$.
\begin{soluzione}
Si ha $f'(x)=e^{x}-e^{-x}-2x$, $f''(x)=e^{x}+e^{-x}-2$, $f'''(x)=e^{x}-e^{-x}$,
$f^{(4)}(x)=e^{x}+e^{-x}$; quindi $f'(0)=f''(0)=f'''(0)=0$ e $f^{(4)}(0)=2>0$. Per il
\cref{teo:criterio-ordine-n} con $n=4$ (pari), l'origine è di minimo relativo stretto.
In accordo, dalla \eqref{eq:taylor-exp} si ricava
$e^{x}+e^{-x}=2+x^{2}+\frac{x^{4}}{12}+o(x^{4})$, cioè $f(x)=\frac{x^{4}}{12}+o(x^{4})$.
\emph{Scelta logica:} il criterio della derivata seconda non decide, perché
$f''(0)=0$; si prosegue fino alla prima derivata non nulla, la cui parità decide.
\end{soluzione}
\end{esercizio}

\begin{esercizio}
Calcolare un valore approssimato di $\sqrt{1{,}1}$ con il polinomio di Taylor di ordine
$2$ di $\sqrt{1+x}$, e stimare l'errore.
\begin{soluzione}
Dalla \eqref{eq:taylor-potenza} con $b=\frac12$:
$\sqrt{1+x}=1+\frac x2-\frac{x^{2}}8+o(x^{2})$. Per $x=0{,}1$ si ottiene
$1+0{,}05-0{,}00125=1{,}04875$, mentre $\sqrt{1{,}1}=1{,}048809\dots$: l'errore è circa
$5{,}9\cdot10^{-5}$. Il termine successivo dello sviluppo è
$\frac{\frac12(-\frac12)(-\frac32)}{3!}x^{3}=\frac{x^{3}}{16}=6{,}25\cdot10^{-5}$, che dà
l'ordine di grandezza corretto dell'errore. \emph{Scelta logica:} la formula di Taylor
con il resto di Peano dice che l'errore è piccolo rispetto a $x^{2}$, ma non quanto; una
stima effettiva si ottiene guardando il primo termine trascurato (e, in modo rigoroso,
con la forma di Lagrange del resto, \cref{teo:resto-lagrange}).
\end{soluzione}
\end{esercizio}

\needspace{8\baselineskip}
\subsection*{Studi di funzione}

\begin{esercizio}\label{es:studio-razionale}
Studiare la funzione $f(x)=\dfrac{x^{2}}{x-1}$ e disegnarne il grafico.
\begin{soluzione}
\textbf{A.} Dominio: $x\neq1$. \textbf{B.} Nessuna simmetria; $f(x)=0$ solo per $x=0$;
$f>0$ per $x>1$, $f\le0$ per $x<1$. \textbf{C.} $f(x)\to\pm\infty$ per $x\to\pm\infty$ e
per $x\to1^{\pm}$: la retta $x=1$ è asintoto verticale. \textbf{D.} Si ha
\[
  f'(x)=\frac{2x(x-1)-x^{2}}{(x-1)^{2}}=\frac{x(x-2)}{(x-1)^{2}} ,
\]
positiva per $x<0$ e per $x>2$, negativa in $(0,1)$ e in $(1,2)$: $x=0$ è di massimo
relativo, con $f(0)=0$, e $x=2$ è di minimo relativo, con $f(2)=4$. \textbf{E.}
Scrivendo $f(x)=x+1+\frac1{x-1}$ (divisione tra polinomi) si ricava facilmente
$f''(x)=\frac{2}{(x-1)^{3}}$: $f$ è concava per $x<1$ e convessa per $x>1$; non ci sono
flessi. \textbf{F.} Dalla stessa scrittura, $f(x)-(x+1)=\frac1{x-1}\to0$ per
$x\to\pm\infty$: la retta $y=x+1$ è asintoto obliquo da entrambe le parti (con la
\cref{prop:asintoto-obliquo}: $m=\lim f(x)/x=\lim x/(x-1)=1$ e
$q=\lim\bigl(f(x)-x\bigr)=\lim x/(x-1)=1$). Il grafico è nella
\cref{fig:studio-razionale}. Si noti che il massimo relativo $f(0)=0$ è minore del minimo
relativo $f(2)=4$: «massimo» e «minimo» relativi sono nozioni locali.
\emph{Scelta logica:} la divisione tra polinomi semplifica tutti i passi successivi:
rende evidenti l'asintoto obliquo, la posizione del grafico rispetto a esso (sopra per
$x>1$, sotto per $x<1$) e il calcolo di $f''$.
\end{soluzione}
\end{esercizio}

\begin{figure}[htb]
\centering
\begin{tikzpicture}
\begin{axis}[width=0.8\textwidth, height=7cm, xmin=-3.7, xmax=5.7, ymin=-5, ymax=8,
  xtick={2}, ytick={4}, xlabel={$x$}, ylabel={$y$}, clip=false]
  \addplot[tratteggio, domain=-3.5:5.5, samples=2] {x+1};
  \addplot[tratteggio] coordinates {(1,-5) (1,8)};
  \addplot[curva, domain=-3.5:0.84, samples=150] {x^2/(x-1)};
  \addplot[curva, domain=1.2:5.5, samples=150] {x^2/(x-1)};
  \addplot[tratteggio] coordinates {(2,0) (2,4) (0,4)};
  \addplot[only marks, mark=*, mark size=1.6pt, rame] coordinates {(0,0) (2,4)};
  \node[font=\scriptsize, grigio, anchor=west] at (axis cs:4.6,4.9) {$y=x+1$};
  \node[font=\scriptsize, blu, anchor=west] at (axis cs:1.25,7.4) {$y=\frac{x^{2}}{x-1}$};
\end{axis}
\end{tikzpicture}
\caption{Il grafico di $f(x)=\frac{x^{2}}{x-1}$ (\cref{es:studio-razionale}): asintoto
verticale $x=1$, asintoto obliquo $y=x+1$, massimo relativo in $(0,0)$ e minimo relativo
in $(2,4)$.}
\label{fig:studio-razionale}
\end{figure}

\begin{esercizio}\label{es:e-pi}
Studiare la funzione $f(x)=\dfrac{\log x}{x}$ e dedurne quale dei due numeri $e^{\pi}$ e
$\pi^{e}$ sia il maggiore.
\begin{soluzione}
Il dominio è $(0,+\infty)$. Si ha $f(x)\to-\infty$ per $x\to0^{+}$ (asintoto verticale
$x=0$) e $f(x)\to0$ per $x\to+\infty$ per la \eqref{eq:hopital-log-potenza} (asintoto
orizzontale $y=0$). La derivata $f'(x)=\frac{1-\log x}{x^{2}}$ è positiva per $x<e$ e
negativa per $x>e$: $f$ è strettamente crescente in $(0,e]$ e strettamente decrescente in
$[e,+\infty)$, con massimo assoluto $f(e)=\frac1e$. La derivata seconda
$f''(x)=\frac{2\log x-3}{x^{3}}$ cambia segno in $x=e^{3/2}$, che è un punto di flesso.
Poiché $\pi>e$ e $f$ è strettamente decrescente in $[e,+\infty)$, risulta
$\frac{\log\pi}{\pi}<\frac{\log e}{e}=\frac1e$, cioè $e\log\pi<\pi$, e quindi, per la
monotonia dell'esponenziale, $\pi^{e}=e^{e\log\pi}<e^{\pi}$. Numericamente,
$\pi^{e}\approx22{,}46$ ed $e^{\pi}\approx23{,}14$. \emph{Scelta logica:} il confronto tra
$a^{b}$ e $b^{a}$ si trasforma, prendendo i logaritmi e dividendo per $ab$, nel confronto
tra $\frac{\log a}a$ e $\frac{\log b}b$, cioè tra due valori della stessa funzione: lo
studio della monotonia risolve il problema per tutte le coppie di numeri distinti non
minori di $e$.
\end{soluzione}
\end{esercizio}
